\part[Гармонический анализ]{Гармонический анализ}

\chapter{Достаточные условия сходимости тригонометрического ряда Фурье в точке.}\label{chapter17}

\section{Пространство интегрируемых функций. Лемма Римана}

\begin{defn}
Функция $f$ называется \textit{абсолютно интегрируемой\rindex{функция!абсолютно интегрируемая}} на промежутке $I \subset \bbR$, если интеграл (вообще говоря, несобственный с конечным числом особенностей) от функции $f$ по промежутку $I$ абсолютно сходится. Множество функций, абсолютно интегрируемых на промежутке $I$, обозначается $L_R(I)$ или $L_R^1(I)$.
\end{defn}

\begin{notion}
$\int\limits_I f(x)\,dx$ сходится абсолютно, если $f$ интегрируема по Риману на любом конечном отрезке $[a; b] \subset I$, не содержащем особенностей, и $\int\limits_I |f(x)|\,dx$ сходится (иначе говоря, $\int\limits_I f(x)\,dx$ и $\int\limits_I |f(x)|\,dx$ сходятся; см. определения 13.2 и 13.3 и теорему 13.10). Одной сходимости $\int\limits_I |f(x)|\,dx$ мало. Например, функция
$$
f(x) = \begin{cases} 1, & \text{если } x \in \bbQ; \\ -1, & \text{если } x \notin \bbQ \end{cases}
$$
не интегрируема по Риману ни на каком конечном отрезке (аналогично примеру 12.1), а модуль её интегрируем.
\end{notion}

Легко видеть, что $L_R(I)$ образует линейное пространство относительно обычных операций сложения и умножения на действительное число (теоремы 13.1 и 13.16).

\begin{defn}
Функция $f$ называется \textit{абсолютно интегрируемой с квадратом\rindex{функция!абсолютно интегрируемая с квадратом}} на промежутке $I \subset \bbR$, если $f \in L_R(I)$ и $f^2 \in L_R(I)$ (т.е. $\int\limits_I f(x)\,dx$ и $\int\limits_I (f(x))^2\,dx$ сходятся абсолютно). Множество таких функций обозначается $L_R^2(I)$.
\end{defn}

\begin{defn}
Линейное пространство $L$ называется \textit{нормированным\rindex{пространство!нормированное}} (\textbf{ЛНП} "--- линейное нормированное пространство), если на нём определена действительнозначная функция $\|\cdot\|$ (\textit{норма\rindex{норма}}) такая, что для любых $x, y \in L$, $\alpha \in \bbR$ имеют место соотношения:
\begin{enumerate}
\item[$1^\circ$.] $\|x + y\| \le \|x\| + \|y\|$ (неравенство треугольника);
\item[$2^\circ$.] $\|\alpha x\| = |\alpha| \cdot \|x\|$;
\item[$3^\circ$.] $\|x\| \ge 0$, причём $\|x\| > 0$ при $x \neq 0$.
\end{enumerate}
\end{defn}

В евклидовом пространстве эти соотношения выполняются для функции $\|x\| = \sqrt{(x, x)}$, поэтому пространство $L_R^2(I)$ является ЛНП с нормой $\|f\|_2 = \sqrt{\int\limits_I (f(x))^2\,dx}$. В пространстве $L_R^1(I)$ нельзя ввести скалярное произведение, но можно ввести норму $\|f\|_1 = \int\limits_I |f(x)|\,dx$. Пространство $L_R^2(I)$ является подмножеством линейного пространства $L_R^1(I)$, поэтому для всех функций из $L_R^2(I)$ определены $\|f\|_2$ и $\|f\|_1$.

\begin{thm}[лемма Римана об осцилляции]
Если $f \in L_R(I)$, где $I$ "--- произвольный промежуток в $\bbR$, то
$$
\lim_{t \to \infty} \int\limits_I f(x) \cos tx\,dx = 0, \quad \lim_{t \to \infty} \int\limits_I f(x) \sin tx\,dx = 0.
$$
\end{thm}


\section{Неравенство Бесселя и равенство Парсеваля}

\begin{thm}
Пусть $e_1, e_2, \ldots, e_n, \ldots$ "--- счётная ОНС в бесконечномерном евклидовом пространстве $L$; $x \in L$; $c_1, c_2, \ldots, c_n, \ldots$ "--- коэффициенты Фурье $x$ по данной ОНС. Тогда ряд $\sum\limits_{k=1}^{\infty} c_k^2$ сходится и
$$
\sum\limits_{k=1}^{\infty} c_k^2 \le \|x\|^2 \quad \text{(\textit{неравенство Бесселя\rindex{неравенство Бесселя}})}.
$$
\end{thm}

\begin{proof}
Из леммы 22.8 следует, что $\sum\limits_{k=1}^n c_k^2 \le \|x\|^2$ при всех $n = 1, 2, \ldots$ Отсюда и следует утверждение теоремы.
\end{proof}

\begin{defn}
Пусть $e_1, e_2, \ldots, e_n, \ldots$ "--- счётная ОНС в бесконечномерном евклидовом пространстве $L$. Говорят, что для элемента $x \in L$ выполняется \textit{равенство Парсеваля\rindex{равенство Парсеваля}} по системе $e_1, e_2, \ldots, e_n, \ldots$, если неравенство Бесселя обращается в равенство
$$
\sum\limits_{k=1}^{\infty} c_k^2 = \|x\|^2.
$$
\end{defn}




\section{Интегральное представление частичных сумм рядов Фурье по тригонометрической системе. Ядро Дирихле}

Рассмотрим $n$-ю частичную сумму ряда Фурье функции $f(x)\in \accentset{*}{L}^{R}_1(-\pi;\pi)$ по тригонометрической системе, ортогональной на интервале $(-\pi;\pi)$. Как показано в прошлом пункте, имеют место формулы~\eqref{2018:ch17:eq18} и  \eqref{2018:ch17:eq19}. Подставляя одну в другую, получим
$$
T_{n}(f;x) = \frac{1}{2\pi} \int_{-\pi}^{\pi} f(\xi) \sum_{\nu = -n}^{n} e^{i\nu(x-\xi)}\,d\xi
$$
\begin{defn}
Функция $D_n(x) = \sum_{\nu = -n}^{n} e^{i\nu x}$ называется \textit{ядром Дирихле порядка $n$}.
\end{defn}

Очевидно, 
$$
D_n(x) = 1 + 2 \sum_{\nu = -n}^{n} \cos{\nu x}
$$
Отсюда следует, что $D_n(x)$ "--- четная $2\pi$-периодическая функция, $D_n(0) = 1+2n$ и
$$
\frac{1}{\pi}\int_{0}^{\pi} D_n(x) \,dx =1
$$
По формуле для суммы геометрической прогрессии получаем, что
$$
D_n(x) = \frac{e^{-inx} - e^{-inx +ix}}{1-e^{ix}} = 
\frac{\sin(n+\frac{1}{2})x}{\sin{\frac{x}{2}}}
$$
для любого $x\ne 2k\pi$, $k \in \bbZ$ (и равно $2n+1$ в противоположном случае, как уже было замечено).

С помощью ядра Дирихле $n$-я частичная сумма ряда Фурье функции $f(x)$ выражается следующим образом:
$$
T_{n}(f;x) = \frac{1}{2\pi} \int_{-\pi}^{\pi} f(\xi) D_n(\xi-x)\,d\xi = \frac{1}{2\pi} \int_{-\pi-x}^{\pi-x} f(\xi + x) D_n(\xi)\,d\xi.
$$
А так как подынтегральная функция является $2\pi$-периодической, то 
$$
T_{n}(f;x)= \frac{1}{2\pi} \int_{-\pi}^{\pi} f(\xi + x) D_n(\xi)\,d\xi.
$$


\section{Тригонометрические ряды Фурье и их сходимость}

\begin{defn}
Пусть $f \in L_R[-l, l]$. Тогда числа
\begin{equation} \label{FourierCoef}
\begin{aligned}
a_n &= \frac{1}{l} \int\limits_{-l}^l f(x) \cos \frac{\pi n x}{l}\,dx, \quad n = 0, 1, 2, \ldots, \\
b_n &= \frac{1}{l} \int\limits_{-l}^l f(x) \sin \frac{\pi n x}{l}\,dx, \quad n = 1, 2, \ldots
\end{aligned}
\end{equation}
называются \textit{коэффициентами Фурье\rindex{коэффициенты Фурье}} функции $f$, а соответствующий функциональный ряд "--- \textit{рядом Фурье\rindex{ряд Фурье}} функции $f$.
\end{defn}

Интегралы, входящие в выражения для коэффициентов Фурье $a_n$ и $b_n$, абсолютно сходятся для всех $f \in L_R[-l, l]$, а не только для $f \in L_R^2[-l, l]$. Абсолютная сходимость их следует из неравенств $\left| f(x) \cos \frac{\pi n x}{l} \right| \le |f(x)|$, $\left| f(x) \sin \frac{\pi n x}{l} \right| \le |f(x)|$ и признака сравнения сходимости несобственных интегралов от знакопостоянных функций.

Из леммы Римана следует, что коэффициенты Фурье абсолютно интегрируемой функции стремятся к нулю при $n \to \infty$. Функцию $f \in L_R[-l, l]$ можно продолжить на всю числовую прямую с периодом $2l$; значения интегралов в формулах \eqref{FourierCoef} не изменятся, если интеграл рассматривать по любому отрезку длины $2l$. Поэтому можно говорить о ряде Фурье функции $f \in L_R[a, a + 2l]$ с периодом $2l$.

Частичная сумма ряда Фурье имеет вид
$$
S_n(f, x) = \frac{a_0}{2} + \sum_{k=1}^n \left( a_k \cos \frac{\pi k x}{l} + b_k \sin \frac{\pi k x}{l} \right).
$$

\begin{thm}[признак Дини сходимости рядов Фурье]
Пусть функция $f \in L_R[-l, l]$ и имеет период $2l$, а $S_0$ "--- такое число, что при некотором $\delta > 0$ сходится интеграл
$$
\int\limits_{0\leftarrow}^{\delta} \frac{|f(x_0 + z) + f(x_0 - z) - 2 S_0|}{z}\,dz.
$$
Тогда ряд Фурье функции $f$ сходится в точке $x_0$ к числу $S_0$.
\end{thm}

\begin{proof}
По лемме 22.7 функция $f$ абсолютно интегрируема на любом конечном отрезке. Поэтому расходимость интеграла могла бы произойти только благодаря особенности в точке $0$, но по условию теоремы этот интеграл сходится; значит, сходится интеграл от такой же функции и на $(0; l]$.

Так как $\int\limits_{-\pi}^{\pi} D_n(u)\,du = \int\limits_{-\pi}^{\pi} \left( \frac{1}{2} + \cos u + \ldots + \cos nu \right) du = \pi$, то в силу чётности функции $\int\limits_0^{\pi} D_n(u)\,du = \frac{\pi}{2}$ и $\int\limits_0^l D_n\!\left( \frac{\pi z}{l} \right) dz = \frac{l}{2}$ (сделана замена $u = \frac{\pi z}{l}$ в интеграле от непрерывной функции). Поэтому $S_0 = \frac{2}{l} \int\limits_0^l S_0 \cdot D_n\!\left( \frac{\pi z}{l} \right) dz$. В силу представления частичной суммы через ядро Дирихле
\begin{equation} \label{FourierDiff}
S_n(f, x_0) - S_0 = \frac{1}{l} \int\limits_0^l \frac{f(x_0 + z) + f(x_0 - z) - 2 S_0}{2 \sin \frac{\pi z}{2l}} \sin\!\left( n + \tfrac{1}{2} \right)\!\frac{\pi z}{l}\,dz.
\end{equation}
Но
$$
\frac{f(x_0 + z) + f(x_0 - z) - 2 S_0}{2 \sin \frac{\pi z}{2l}} = \frac{f(x_0 + z) + f(x_0 - z) - 2 S_0}{z} \cdot \varphi(z),
$$
где $\varphi(z) = \dfrac{z}{2 \sin \frac{\pi z}{2l}}$. На отрезке $[0; l]$ знаменатель обращается в нуль только при $z = 0$, а $\lim\limits_{z \to 0} \varphi(z) = \frac{l}{\pi}$, поэтому функция $\varphi(z)$ ограничена на $(0; l]$. Так как $\frac{f(x_0 + z) + f(x_0 - z) - 2 S_0}{z} \in L_R[0, l]$, то и $\frac{f(x_0 + z) + f(x_0 - z) - 2 S_0}{2 \sin \frac{\pi z}{2l}} \in L_R[0, l]$. По лемме Римана интеграл в правой части \eqref{FourierDiff} стремится к нулю при $n \to \infty$, поэтому $\lim\limits_{n \to \infty} S_n(f, x_0) = S_0$.
\end{proof}

\begin{defn}
Пусть функция $f$ определена в некоторой $\dnei{x_0}$ и существуют конечные односторонние пределы $f(x_0 + 0)$ и $f(x_0 - 0)$, причём
$$
\exists\, \alpha > 0,\ \exists\, C > 0 : \forall\, z \in (0; \delta) \longrightarrow
$$
$$
|f(x_0 + z) - f(x_0 + 0)| \le C z^\alpha, \quad |f(x_0 - z) - f(x_0 - 0)| \le C z^\alpha.
$$
Тогда говорят, что функция $f$ удовлетворяет в точке $x_0$ \textit{условию Липшица\rindex{условие Липшица}} (или \textit{условию Гёльдера\rindex{условие Гёльдера}}) порядка $\alpha$. Множество таких функций обозначается $\Lip_\alpha(x_0)$.
\end{defn}

\begin{notion}
Ясно, что при $\alpha > \beta > 0$ имеет место включение $\Lip_\alpha(x_0) \subset \Lip_\beta(x_0)$.
\end{notion}

\begin{thm}[признак Липшица сходимости рядов Фурье]
Пусть функция $f \in L_R[-l, l]$ и имеет период $2l$, причём $f \in \Lip_\alpha(x_0)$, $\alpha > 0$. Тогда ряд Фурье функции $f$ сходится в точке $x_0$ к значению $S_0 = \dfrac{f(x_0 + 0) + f(x_0 - 0)}{2}$.
\end{thm}

\begin{proof}
При всех $z \in (0; \delta)$
\begin{multline*}
\frac{|f(x_0 + z) + f(x_0 - z) - 2 S_0|}{z} \le \\
\le \frac{|f(x_0 + z) - f(x_0 + 0)| + |f(x_0 - z) - f(x_0 - 0)|}{z} \le \frac{C z^\alpha + C z^\alpha}{z} = 2C \cdot \frac{1}{z^{1 - \alpha}}.
\end{multline*}
Так как $\alpha > 0$, то $1 - \alpha < 1$, и $\int\limits_{0\leftarrow}^{\delta} \frac{dz}{z^{1 - \alpha}}$ сходится. Тогда $\int\limits_{0\leftarrow}^{\delta} \frac{|f(x_0 + z) + f(x_0 - z) - 2 S_0|}{z}\,dz$ сходится по признаку сравнения. По признаку Дини ряд Фурье функции $f$ сходится в точке $x_0$ к значению $S_0$.
\end{proof}

\begin{defn}
Говорят, что функция $f$, определённая в некоторой проколотой правой (левой) окрестности точки $x_0$, имеет \textit{обобщённую правую (левую) одностороннюю производную\rindex{производная!обобщённая односторонняя}}, если существуют конечные пределы $f(x_0 + 0)$ и $\lim\limits_{t \to +0} \frac{f(x_0 + t) - f(x_0 + 0)}{t}$ (соответственно $f(x_0 - 0)$ и $\lim\limits_{t \to +0} \frac{f(x_0 - t) - f(x_0 - 0)}{-t}$).
\end{defn}

Геометрический смысл такой обобщённой односторонней производной "--- угловой коэффициент односторонней касательной в точке с абсциссой $x_0$ графика функции $y = f(x)$, доопределённой в этой точке соответствующим предельным значением функции (см. рис. 22.2).

\begin{cons}
Пусть функция $f \in L_R[-l, l]$ имеет период $2l$, причём $f$ непрерывна и имеет конечные односторонние производные $f'_+(x_0)$ и $f'_-(x_0)$ в точке $x_0$. Тогда ряд Фурье $f$ сходится в точке $x_0$ к значению $f(x_0)$.
\end{cons}

\begin{proof}
$f'_+(x_0) = \lim\limits_{t \to +0} \dfrac{f(x_0 + t) - f(x_0)}{t} = A$. В этом случае
$$
\exists\, \delta_1 > 0 : \forall\, t \in (0; \delta_1) \longrightarrow |f(x_0 + t) - f(x_0)| \le (|A| + 1) t.
$$
Аналогично из существования $f'_-(x_0) = B$ следует, что
$$
\exists\, \delta_2 > 0 : \forall\, t \in (0; \delta_2) \longrightarrow |f(x_0 - t) - f(x_0)| \le (|B| + 1) t.
$$
Тогда, если $\delta = \min(\delta_1, \delta_2) > 0$, $C = \max(|A| + 1, |B| + 1)$, то
$$
\forall\, t \in (0; \delta) \longrightarrow |f(x_0 + t) - f(x_0)| \le C t, \quad |f(x_0 - t) - f(x_0)| \le C t.
$$
Так как $f(x_0 + 0) = f(x_0 - 0) = f(x_0)$, то отсюда следует, что $f \in \Lip_1(x_0)$ и ряд Фурье $f$ сходится в точке $x_0$ к значению $\dfrac{f(x_0 + 0) + f(x_0 - 0)}{2} = f(x_0)$.
\end{proof}

\begin{cons}
Пусть функция $f \in L_R[-l, l]$ и имеет период $2l$, причём в точке $x_0$ у неё разрыв первого рода и существуют конечные обобщённые односторонние производные. Тогда ряд Фурье $f$ сходится в точке $x_0$ к значению $\dfrac{f(x_0 + 0) + f(x_0 - 0)}{2}$.
\end{cons}

\begin{proof}
Рассуждение аналогично доказательству следствия 1.
\end{proof}

Следствия 1 и 2 из признака Липшица дают работающие на практике условия сходимости ряда Фурье в точке.

Отметим, что если $f \in L_R[-l, l]$ и чётна, т.е. $f(-x) = f(x)$ на $[-l; l]$, то $b_n = \frac{1}{l} \int\limits_{-l}^l f(x) \sin \frac{\pi n x}{l}\,dx = 0$, $n = 1, 2, \ldots$, т.е. в ряду Фурье есть свободный член и косинусы, но нет синусов; если $f$ нечётна, т.е. $f(-x) = -f(x)$ на $[-l; l]$, то $a_n = \frac{1}{l} \int\limits_{-l}^l f(x) \cos \frac{\pi n x}{l}\,dx = 0$, $n = 0, 1, 2, \ldots$, т.е. в ряду Фурье только синусы.


